%%
%% spring-lecture02.tex
%% 
%% Made by Alex Nelson <pqnelson@gmail.com>
%% Login   <alex@lisp>
%% 
%% Started on  2026-04-04T09:28:20-0700
%% Last update 2026-04-04T09:28:20-0700
%% 

\lecture{}

\begin{definition}
Any $R$-algebra $A$ is the $A$-bimodule by its product.
\begin{enumerate}
\item The \define{Left Regular Representation} of $A$ is $A$ and $\lambda\colon A\to\End_{A}(A)$.
\item The \define{Right Regular Representation} of $A$ is the
  $A$-module $A$ and $\rho\colon A\to\End_{A}(A)$
\end{enumerate}
\end{definition}

\begin{node}
Note: $\lambda$ and $\rho$ are bijective.
Let us check this for $\lambda$. We see $\ker(\lambda)=\{x\in A\mid xA=0\}=\{0\}$
since $1_{A}\in A$ which means for any $x\in\ker(\lambda)$ we have $x1_{A}=x=0$.
\end{node}

\begin{node}
Any $A$-algebra endomorphism $\phi\in\End_{A}(A)$ is determined by
$\phi(1)$. Specifically, for any $x\in A$ we have $\phi(x)=\phi(1_{A}x)=\phi(1_{A})x=\lambda_{\phi(1)}(x)$,
so $\phi=\lambda_{\phi(1)}$.
\end{node}

\subsection{Matrix Algebras}

\begin{definition}
Let $A$ be an $R$-algebra, let $n\in\NN$. The \define{$n\times n$
Matrix Algebra} over $R$ denoted $M_{n}(A)$ consists of all $n\times n$
matrices with entries in $A$ equipped with matrix addition and matrix
multiplication. Scalar multiplication is componentwise multiplication
by $R$.
\end{definition}

\begin{notation}
We write $\canMat_{i,j}$ for the \define{Canonical Matrix} which has
its only nonzero entry located in row $i$ and column $j$ where it is
equal to $1$, and zero at all other entries. Then $\canMat_{i,j}\in M_{n}(A)$
and also
\begin{equation}
\canMat_{i,j}\canMat_{k,\ell}=\delta_{j,k}\delta_{i,\ell}.
\end{equation}
\end{notation}

\begin{notation}
We also write $M_{m,n}(A)$ for the set of $m\times n$ matrices over $A$.
This \emph{is not} an algebra over $A$ \emph{unless} $m=n$. But
$M_{m,n}(A)$ is always an $A$-bimodule. If $\alpha=(x_{i,j})\in M_{m,n}(A)$
and $y\in A$, then $y\alpha=(yx_{i,j})\in M_{m,n}(A)$. Similarly,
$\alpha y\in M_{m,n}(A)$.

We identify $M_{n,n}(A)=M_{n}(A)$ is the free $A$-bimodule with basis
$\{\canMat_{i,j}\mid 1\leq i,j\leq n\}$. Also observe $(\alpha y)\beta=\alpha(y\beta)$
for any $\alpha,\beta\in M_{n}(A)$ and $y\in A$.
\end{notation}

\begin{caution}
If $A$ is not commutative, then $M_{n}(A)$ is not an $A$-algebra.
\end{caution}

\begin{definition}
A \define{Division Algebra} over $R$ consists of an $R$-algebra $D$
such that for any nonzero $x\in D$ there exists an $x^{-1}\in D$ such
that $x\cdot x^{-1}=x^{-1}x=1_{D}$. 
\end{definition}

\begin{definition}
Let $A$ be an $R$-algebra. We say $A$ is \define{Simple} if it has no
nontrivial two-sided ideals.
\end{definition}

\begin{lemma}
If $D$ is a division algebra, then $M_{n}(D)$ is simple.
\end{lemma}

\begin{proof}
Let $I\neq0$ be a two-sided ideal of $M_{n}(D)$. Then there exists
$\beta\in I$ nonzero $\beta\neq0$ with entries $\beta=(y_{i,j})$. So
there is a nonzero entry $y_{r,s}\neq0$. Then for any
$\alpha=(x_{i,j})\in M_{n}(D)$ we have
\begin{subequations}
  \begin{align}
\alpha&=\sum_{i,j}\canMat_{i,j}x_{i,j}\\
&=\sum_{i,j}(\canMat_{i,r}\beta\canMat_{s,j})y_{r,s}^{-1}x_{i,j}\\
&=\sum_{i,j}(y_{r,s}\canMat_{i,j})y_{r,s}^{-1}x_{i,j}
  \end{align}
\end{subequations}
Therefore $\alpha\in I$. But since $\alpha$ was arbitrary, this means
$I=M_{n}(D)$. 
\end{proof}

\begin{xca}[Pierce]
Prove this lemma generalizes from division algebras $A$ to any simple
algebra $A$.
\end{xca}

\begin{example}
Let $R=\FF$ be a field, and $A$ be a vector space over $\FF$ with
basis $x_{1}$, \dots, $x_{n}$. Then a (possibly non-associative
non-unital) $\FF$-algebra structure on $A$ is definedby specifying
\begin{equation}
x_{i}x_{j}=\sum^{n}_{k=1}x_{k}{a_{i,j}}^{k}
\end{equation}
where (for all $i,j,k=1,\dots,n$) the coefficients
${a_{i,j}}^{k}\in\FF$ are called the 
\define{Structure Constants}. We can express properties of the algebra
in terms of these structure constants.
\begin{enumerate}
\item\textsc{Associativity}: if for any $i$, $j$, $k$ we have $(x_{i}x_{j})x_{k}=x_{i}(x_{j}x_{k})$
  then
  \begin{equation}
\sum_{r}{a_{i,j}}^{r}{a_{r,k}}^{s}=\sum_{r}{a_{j,k}}^{r}{a_{i,r}}^{s}
  \end{equation}
  for any $i$, $j$, $k$, $s$.
\item\textsc{Unital}: (possibly reindexing to make $x_{1}=1_{A}$) for
  any $j=1,\dots,n$ we have $x_{1}x_{j}=x_{j}x_{1}=x_{j}$. This means
  $\lambda_{x_{1}}=\rho_{x_{1}}=\id_{A}$ so ${a_{1,j}}^{k}={a_{j,1}}^{k}=\delta_{j,k}$
  for any $j,k=1,\dots,n$.
\item\textsc{Morphisms in terms of structure constants}. Let
  $(A,{a_{i,j}}^{k})$ and $(A,{b_{i,j}}^{k})$ be two algebras with the
  same underlying vector space. The following are equivalent:
\begin{enumerate}
\item They are isomorphic;
\item there exists a $\gamma\in\Aut_{\FF}(A)$ such that
\begin{equation}
\gamma(x_{i})=\sum_{r}x_{r}{c_{i}}^{r}\gamma(x_{i}x_{j})
\end{equation}
\item for all $i$, $j$, $k$ we have
  \begin{equation}
\sum_{r}{c_{r}}^{k}{b_{i,j}}^{r}=\sum_{s,t}{a_{s,t}}^{k}{c_{i}}^{s}{c_{j}}^{t}.
  \end{equation}
\end{enumerate}
\end{enumerate}
\end{example}

\begin{definition}
Let $\FF$ be a field such that $\Char(\FF)\neq2$. Let $a,b\in\FF$ be
fixed nonzero elements of $\FF$. Then the \define{(Generalized) Quaternion Algebra}
is the 4-dimensionalalgebra over $\FF$ denoted $\generalizedquat{a}{b}{\FF}$
with the basis $1$, $\quat{i}$, $\quat{j}$, $\quat{k}$ and product
\begin{equation}
\quat{i}^{2}=a,\quad\quat{j}^{2}=b,\quad\quat{i}\quat{j}=-\quat{j}\quat{i}=\quat{k}
\end{equation}
(and we can then infer $\quat{k}^{2}=-ab$, $\quat{i}\quat{k}=-\quat{k}\quat{i}=\quat{j}a$, and $\quat{j}\quat{k}=-\quat{k}\quat{j}=-\quat{i}b$).
\end{definition}

\begin{example}
The standard quaternions are a special case of the generalized
quaternions $\HH=\generalizedquat{-1}{-1}{\RR}$.
\end{example}

\begin{definition}
Let $A$ be an $R$-algebra. 
\begin{enumerate}
\item The \define{Center} of $A$ is the $R$-subalgebra of $A$ denoted
  $Z(A)=\{z\in A\mid\forall x\in A\ldotp xz=zx\}$.
\item If $R=\FF$ is a field, we say $A$ is \define{Central} if $Z(A)=R$.
  (There is a complicated generalization to the setting of commutative
  rings $R$ which are not necessarily fields.)
\end{enumerate}
\end{definition}

\begin{lemma}\label{lemma:1.10}
For any $A=\generalizedquat{a}{b}{\FF}$, $A$ is central and simple.
\end{lemma}

\begin{proof}
Let $x=c_{0}+\quat{i}c_{1}+\quat{j}c_{2}+\quat{k}c_{3}\in A$ where $c_{0},\dots,c_{3}\in\FF$.
Let us also introduce the notation
\begin{equation}
[x,y]:=xy-yx
\end{equation}
for the commutator of any $x,y\in A$.

\textsc{Claim 1} (central): $x\in Z(A)$ if and only if $c_{1}=c_{2}=c_{3}=0$.
We see
\begin{subequations}
  \begin{align}
[\quat{i},x] &=\quat{j}(2ac_{3})+\quat{k}(2c_{2})\\
[\quat{j},x] &=\quat{i}(2ac_{3})+\quat{k}(-2c_{1})\\
[\quat{k},x] &=\quat{j}(2ac_{2})+\quat{k}(-2c_{1}),
  \end{align}
\end{subequations}
and then $x\in Z(A)$ if and only if the right-hand side of these
commutators vanish (if and only if $c_{1}=c_{2}=c_{3}=0$ as desired).

\textsc{Claim 2} (simple): Let $I$ be an ideal of $A$ and $0\neq x\in I$.
We will show $I=A$. Then
\begin{subequations}
  \begin{align}
[\quat{j},[\quat{i},x]] &=\quat{i}(-4bc_{2})\\
[\quat{k},[\quat{j},x]] &=\quat{j}(4abc_{3})\\
[\quat{i},[\quat{k},x]] &=\quat{i}(-4ac_{1}).
  \end{align}
\end{subequations}
If there exists a nonzero $c_{i}\neq0$ for some $i\neq0$, then $I$ contains a unit (so
$I=A$). If all $c_{i}=0$ for $i\neq0$, then $c_{0}=x\in I$ is a unit; i.e., $I=A$.
\end{proof}

\begin{notation}
We denote $A=\generalizedquat{a}{b}{\FF}=\FF\oplus A_{+}$ where
\begin{equation}
A_{+}=\quat{i}\FF\oplus\quat{j}\FF\oplus\quat{k}\FF.
\end{equation}
For any $x\in A$, we can decompose it as $x=c+z$ where $c\in\FF$ and
$z\in A_{+}$.
\end{notation}

\begin{definition}
Using the preceding notation, for any $x=c+z\in A=\generalizedquat{a}{b}{\FF}$,
we define its \define{Conjugate} to be the element $x^{*}:=c-z$ where
$c\in\FF$ and $z\in A_{+}$.
\end{definition}

\begin{note}
For any $x,y\in\generalizedquat{a}{b}{\FF}$ and $d\in\FF$ we have
\begin{subequations}
\begin{equation}
(x+y)^{*}=x^{*}+y^{*},
\end{equation}
and
\begin{equation}
(xy)^{*}=y^{*}x^{*},
\end{equation}
and it is convolutory
\begin{equation}
(x^{*})^{*}=x
\end{equation}
and leaves the ``non-imaginary part'' fixed
\begin{equation}
d^{*}=d.
\end{equation}
\end{subequations}
\end{note}

\begin{definition}
Let $A=\generalizedquat{a}{b}{\FF}$ be the generalized quaternion
algebra over $\FF$. Let $x\in A$. We define the \define{Norm} of $x$
to be the element $\nu(x)$ of $\FF$ equal to
\begin{equation}
\nu(x):=xx^{*}.
\end{equation}
More explicitly, if $x=c_{0}+\quat{i}c_{1}+\quat{j}c_{2}+\quat{k}c_{3}$,
\begin{equation}
\nu(x)=c_{0}^{2}-ac_{1}^{2}-bc_{2}^{2}-abc_{3}^{2}
\end{equation}
and also $\nu(x)=x^{*}x=\nu(x^{*})$.
\end{definition}

\begin{node}
Observe for any $x,y\in A$ we have
\begin{subequations}
  \begin{align}
\nu(xy) &= xy(xy)^{*}\\
&= xy(y^{*}x^{*})\\
&= x(yy^{*})x^{*}\\
&= xx^{*}\nu(y)\\
&= \nu(x)\nu(y).
  \end{align}
\end{subequations}
Also for any $d\in\FF$, we see $\nu(d)=d^{2}$.
\end{node}

\begin{lemma}\label{lemma:1.11}
Let $A=\generalizedquat{a}{b}{\FF}$. Then the following are equivalent:
\begin{enumerate}
\item $A$ is a division algebra
\item If $x\neq0$ is an element of $A$, then $\nu(x)\neq0$
\item If $c_{0}^{2}=ac_{1}^{2}+bc_{2}^{2}$, then $c_{0}=c_{1}=c_{2}=0$.
\end{enumerate}
\end{lemma}

\begin{proof}
$(1)\implies(2)$ If $x\neq0$, then
  $\nu(x)\nu(x^{-1})=\nu(xx^{-1})\nu(1)=1$, which means $\nu(x)\neq0$.

$(2)\implies(1)$ For any nonzero $x\in A$, we have
\begin{subequations}
  \begin{align}
(x^{*}\nu(x)^{-1})x &= xx^{*}\nu(x)^{-1}\\
    &=\nu(x)\nu(x)^{-1}\\
    &=1,
  \end{align}
\end{subequations}
so $x^{-1}=x^{*}\nu(x)^{-1}$.

$(2)\implies(3)$. Suppose $c_{0}^{2}=ac_{1}^{2}+bc_{2}^{2}$ and
assume for contradiction not all of these coefficients are zero, i.e.,
at least one of these $c_{i}\neq0$. Then
$x=c_{0}+\quat{i}c_{1}+\quat{j}c_{2}+\quat{k}0\neq0$ and
$\nu(x)=0$ which is a contradiction.

$(3)\implies(2)$ [Trick: we prove (3) implies the contrapositive of (2)] For
$x=d_{0}+\quat{i}d_{1}+\quat{j}d_{2}+\quat{k}d_{3}\in A$
such that $\nu(x)=0$, we see
\begin{equation}
d_{0}^{2}-bd_{2}^{2}=a(d_{1}^{2}-bd_{3}^{2}).
\end{equation}
Then
\begin{subequations}
  \begin{align}
a(d_{1}^{2}-bd_{3}^{2})^{2}&=(d_{0}^{2}-bd_{2}^{2})(d_{1}^{2}-bd_{3}^{2})\\
&=(d_{0}d_{1}+bd_{2}d_{3})^{2}+b(d_{0}d_{3}+d_{1}d_{2})^{2}.
  \end{align}
\end{subequations}
Writing
\begin{equation}
c_{0}:=(d_{0}d_{1}+bd_{2}d_{3})^{2},\quad\mbox{and}\quad c_{1}:=b(d_{0}d_{3}+d_{1}d_{2})^{2},
\end{equation}
then $c_{0}=ac_{1}^{2}+bc_{2}^{2}$, so $c_{0}=c_{1}=c_{2}=0$ (which
means $d_{1}=d_{3}=0$). Then
\begin{equation}
d_{0}^{2}-bd_{2}^{2}=0,
\end{equation}
which means
\begin{equation}
d_{0}^{2}=a0^{2}+bd_{2}^{2},
\end{equation}
or equivalently $d_{0}=d_{2}=0$. Hence $x=0$.
\end{proof}

\begin{remark}
Observe: this implies $\HH=\generalizedquat{-1}{-1}{\RR}$ forms a
division algebra since $c_{0}^{2}=-c_{1}^{2}-c_{2}^{2}$ implies
$c_{0}=c_{1}=c_{2}=0$ for any $c_{0},c_{1},c_{2}\in\RR$.
\end{remark}